% Auto-generated by scripts/aggregate_section_5_7.py
% Cycle-zero composite calibration (per-signal Pearson rho + boost-layer weights).
% Source: outputs/cycle_0/composite_calibration.json (boost weights, intercept).
% Pearson rho re-computed from outputs/cycle_0/calibration/calibration_fold_samples.json.
\begin{table}[!htbp]
\centering\small
\begin{tabular}{lccr}
\toprule
\textbf{Signal} & \textbf{Pearson $\rho$} & \textbf{Direction} & \textbf{Boost weight} \\
\midrule
$p^{\text{mean}}_{\text{ground}}$ & +0.101 & increasing & +0.184 \\
$p^{\text{max}}_{\text{ground}}$ & +0.140 & increasing & +0.379 \\
$p^{\text{atomic}}_{\text{ground}}$ & +0.101 & increasing & +0.184 \\
$p_{\text{entail}}$ & -0.156 & decreasing & +0.067 \\
$q_{a,\text{rel}}$ & +0.299 & increasing & +0.491 \\
$u_{\text{internal}}$ & +0.425 & increasing & +0.334 \\
$u_{\text{token}}$ & +0.119 & increasing & +0.126 \\
$u_{\text{dropout}}$ & -0.411 & decreasing & +0.190 \\
$s_{\text{avg}}$ & +0.098 & increasing & +0.175 \\
\midrule
\multicolumn{3}{r}{Boost intercept} & +0.265 \\
\bottomrule
\end{tabular}
\caption{Cycle-zero calibrated probability composite under the locked configuration. Each signal is mapped to a calibrated probability through per-signal isotonic regression on the $n = 1500$ training-panel calibration fold ($\mathrm{em\ rate} = 0.446$); the boost layer aggregates the calibrated signals through logistic regression with L2 regularisation strength $C = 0.01$. The Pearson $\rho$ column reports the per-signal correlation with the exact-match label and determines the per-signal isotonic direction. The boost weight column reports the per-signal weight learned by the regularised logistic boost layer. The composite is shared across the three training benchmarks; the per-benchmark child curves are blended toward this pooled fit through the shrinkage prior at $\alpha_{\text{shrink}} = 0.6$ from \Cref{sec:composite-shrinkage}. The signal $h_{\text{norm}}$ is a calibration-time-only diagnostic that does not participate in the deployed composite and is therefore absent from the table.}
\label{tab:composite-calibration}
\end{table}
